% Generated draft: requires mathematical and source review.
\documentclass[11pt]{article}
\usepackage[margin=25mm]{geometry}
\usepackage{fontspec}
\setmainfont{Latin Modern Roman}
\usepackage{amsmath,amssymb,mathtools,unicode-math}
\setmathfont{Latin Modern Math}
\usepackage{enumitem,needspace,xurl,hyperref}
\setlength{\parindent}{0pt}
\setlength{\parskip}{5pt}
\setlength{\emergencystretch}{3em}
\newcommand{\N}{\mathbb{N}}
\newcommand{\Z}{\mathbb{Z}}
\newcommand{\Q}{\mathbb{Q}}
\newcommand{\R}{\mathbb{R}}
\newcommand{\C}{\mathbb{C}}
\newcommand{\F}{\mathbb{F}}
\newcommand{\A}{\mathbb{A}}
\newcommand{\PP}{\mathbb{P}}
\newcommand{\E}{\mathbb{E}}
\newcommand{\eps}{\varepsilon}
\newcommand{\dd}{\,\mathrm{d}}
\begin{document}
\section*{UnsolvedMath ID: 701522}
\textbf{Source problem code:} catalogue entry
\subsection{Definitions and mathematical statement}
Work in the usual finite-group interpretation of formation theory; the supplied record mentions a fixed finite group but does not explicitly state the ambient class of groups. A formation is a nonempty class of finite groups, closed under isomorphisms, homomorphic images, and finite subdirect products. For a formation \(\mathcal H\) and a finite group \(G\), define the \(\mathcal H\)-residual by
\[
G^{\mathcal H}=\bigcap\{N\trianglelefteq G:G/N\in\mathcal H\}.
\]
Use the customary formation product
\[
\mathcal A\mathcal B=\{G:G\text{ is finite and }G^{\mathcal B}\in\mathcal A\}.
\]
These conventions make explicit the standard interpretation of terminology that the supplied record does not define. Write \(\Phi(G)\) for the intersection of the maximal subgroups of \(G\), with \(\Phi(1)=1\). A formation \(\mathcal S\) is saturated if, for every finite group \(G\),
\[
G/\Phi(G)\in\mathcal S\quad\Longrightarrow\quad G\in\mathcal S.
\]
For a fixed finite group \(X\), let
\[
\mathcal M=\bigcap\{\mathcal S:\mathcal S\text{ is a saturated formation and }X\in\mathcal S\}.
\]
Thus \(\mathcal M\) is the one-generated saturated formation specified in the source. Let \(\mathcal F\subseteq\mathcal M\) be a formation satisfying
\[
\mathcal F\ne\mathcal F\mathcal F.
\]
The following questions are asked for every such choice of \(X\) and \(\mathcal F\).
\begin{enumerate}
\item[(a)] Must there exist a positive integer \(t\) and formations \(\mathcal F_1,\ldots,\mathcal F_t\) such that
\[
\mathcal F=\mathcal F_1\cdots\mathcal F_t
\]
and each \(\mathcal F_i\) is non-decomposable?
\item[(b)] If such a factorization into non-decomposable formations exists, are all its factors uniquely determined?
\end{enumerate}
The supplied record does not define ``non-decomposable.'' This term remains an unresolved predicate in both questions; in particular, the admissible decompositions and any restrictions on their factors require confirmation. The record also does not specify whether uniqueness in (b) concerns the ordered sequence of factors or the factors with multiplicities up to permutation. Consequently, (b) retains the source's uniqueness question without imposing either interpretation.
\subsection{Short English statement}
A. N. Skiba asks whether every formation \(\mathcal F\) contained in a one-generated saturated formation and satisfying \(\mathcal F\ne\mathcal F\mathcal F\) admits a finite factorization into non-decomposable formations, and whether the factors of any such factorization are uniquely determined. The meaning of non-decomposability and the intended equivalence of factorizations require source clarification.
\subsection{Sources}
\begin{itemize}
\item Supplied local UnsolvedMath registry record: permanent ID \(701522\), exact source problem code \texttt{catalogue entry}, title ``Kourovka Notebook Problem 18.93.'' The registry labels the problem \emph{open}; this is a supplied status, not an independently verified mathematical or literature assessment.
\item Attributed in the supplied record to A. N. Skiba, proposed in 2014. Source attribution: \emph{The Kourovka Notebook (Unsolved Problems in Group Theory)}, 18th issue (2014), problem 18.93; the record additionally identifies the 21st issue edition (2026), \texttt{arXiv:1401.0300v48}.
\item The supplied collection description identifies E. I. Khukhro and V. D. Mazurov as editors of the Kourovka Notebook collection covering Issues 1--21 (1965--2026). The registry record was updated on \texttt{2026-10-16}; this timestamp is not a supplied literature-verification date. The cited source has not been independently consulted for this draft.
\end{itemize}
\end{document}
